\chapter*{Conclusion générale}
\addcontentsline{toc}{chapter}{Conclusion générale}
\markboth{Conclusion générale}{Conclusion générale}

La problématique de ce projet de fin d'études était de déterminer comment automatiser et coordonner des migrations logicielles Java/Spring Boot et Angular en combinant les capacités de raisonnement des agents avec des mécanismes déterministes, tout en préservant l'intégrité du code source, la traçabilité des décisions, la reproductibilité des opérations et le contrôle du développeur.

Pour répondre à cette problématique, l'objectif a été de concevoir, d'implémenter et d'évaluer une \textbf{Migration Factory intelligente, extensible et gouvernée}. La contribution principale réside dans une architecture multi-agents hybride et centralement orchestrée, dans laquelle les opérations critiques restent sous contrôle déterministe et supervision humaine. Le développeur conserve ainsi l'autorité sur les décisions sensibles tandis que la plateforme structure l'analyse, la transformation, la validation, la réparation et la conservation des preuves.

L'étude du processus manuel a confirmé que la difficulté d'une migration ne réside pas uniquement dans la modification du code. Elle concerne également la préparation des environnements, la compatibilité des versions, l'enchaînement des validations, le diagnostic des échecs et la conservation des décisions. Ces constats ont justifié la séparation entre raisonnement agentique, orchestration persistante, exécution déterministe et points de décision humaine.

L'évolution du projet constitue un premier résultat concret. Une première démonstration Java/Spring Boot, réalisée le 25 mai 2026, a validé la faisabilité d'un workflow en ligne de commande. Une deuxième démonstration, le 3 juin, a présenté une version CLI corrigée et consolidée à une audience CGI plus large. La troisième démonstration, le 19 juin, a marqué le passage à une \textit{Control Tower} full stack intégrant backend, frontend, persistance, cockpit et chatbot.

La déclinaison Java/Spring Boot constitue ainsi le système le plus mature. Elle valide les principaux mécanismes du projet : trajectoires progressives, points de validation, isolation, compilation et tests, réparation gouvernée, cockpit, assistant et reporting. Les travaux engagés depuis le 20 juin préparent une V1 comprenant un nouveau frontend, une boucle de réparation renforcée et une meilleure exploitation des LLM. À la date de référence du rapport, fixée au 30 août 2026, le jalon de démonstration associé à cette V1 était prévu pour septembre 2026.

Le choix d'Azure OpenAI renforce l'intégration de la solution dans le contexte de CGI. La plateforme ne dépend pas d'un modèle unique codé en dur : l'endpoint, l'authentification et les noms de déploiement sont externalisés. CGI peut ainsi connecter les modèles qu'elle a déployés et gouvernés dans sa propre ressource Azure, tout en attribuant des déploiements distincts aux rôles de proposition, de révision, d'assistance et de secours. L'implémentation Java actuelle utilise une authentification par clé API ; l'adoption de Microsoft Entra ID, d'une identité managée et d'un coffre de secrets constitue une évolution d'industrialisation.

La déclinaison Angular montre que les principes peuvent être adaptés à un second écosystème, avec des profils Node.js, npm et TypeScript ainsi que des transformations propres au frontend. Son développement a commencé à la mi-juillet 2026 et l'objectif a ensuite été étendu à une trajectoire progressive d'Angular 11 vers Angular 21. La preuve d'exécution versionnée retenue valide les transitions Angular 18 vers 19 et 19 vers 20 ; le stage 20 vers 21 est préparé mais non démarré. Le parcours 11 vers 21 demeure donc un objectif d'achèvement et non un résultat déjà démontré. L'utilisation effective d'Azure dans Angular devra également être confirmée par l'audit du dépôt correspondant.

Les résultats disponibles montrent que la Factory peut structurer une migration par stages, conserver un état persistant, relier les décisions aux preuves, isoler les modifications et imposer une revalidation avant de poursuivre. Ils confirment également l'intérêt de la séparation entre raisonnement, décision et exécution pour renforcer la supervision et la traçabilité du processus.

Néanmoins, l'évaluation quantitative demeure partielle. Les gains de temps, les taux de réussite, la consommation des modèles et le comportement sur un ensemble représentatif d'applications doivent encore être mesurés de manière reproductible. Une démonstration constitue une validation par les parties prenantes ; elle ne remplace pas une campagne expérimentale répétée et instrumentée.

Les perspectives sont organisées selon quatre axes complémentaires.

\textbf{À court terme}, les travaux portent sur l'achèvement du parcours Angular, la consolidation de la réparation gouvernée, du reporting et de l'assistant, ainsi que sur la stabilisation des scénarios de validation de bout en bout.

\textbf{Pour l'industrialisation}, les priorités concernent l'authentification et les rôles, l'utilisation de Microsoft Entra ID et d'Azure Key Vault, la mise en place d'une chaîne d'intégration continue, l'amélioration de la portabilité et l'adaptation de la persistance à des déploiements multi-instance.

\textbf{Pour l'évaluation}, il reste à automatiser la collecte des durées, des appels LLM, des coûts et des interventions humaines, puis à constituer un corpus reproductible de projets Java et Angular afin d'exécuter plusieurs campagnes comparables.

\textbf{À plus long terme}, les réparations validées pourraient être capitalisées sous forme de règles déterministes réutilisables, tandis que de nouveaux profils et adaptateurs permettraient d'étendre la démarche à d'autres versions, frameworks et écosystèmes.

La progression du prototype CLI vers une plateforme full stack, puis vers une V1 gouvernée et une déclinaison Angular, constitue ainsi une base crédible pour poursuivre l'industrialisation d'une plateforme générale de migration logicielle.
